\documentclass{article}

\usepackage[english]{babel}
\usepackage[letterpaper,top=2cm,bottom=2cm,left=3cm,right=3cm]{geometry}
\usepackage{amsmath,amssymb}
\setlength{\parindent}{0pt}

\title{DSC 190/291: Learning Theory through Formal Proof and Proof Presentation\\Assignment 5}
\date{UCSD, Spring 2026}

\begin{document}
\maketitle

\section{Overview}

This assignment has three parts:
\begin{itemize}
\item Part A: sparse halfspaces and agnostic hardness.
\item Part B: surrogate losses and the fixed-feature barrier.
\item Part C: AI usage report.
\end{itemize}

\section{Part A: Sparse Halfspaces and Agnostic Hardness (40 points)}

Let
$$
\mathcal{H}_{d,k}
=
\left\{
x \mapsto \operatorname{sign}(\langle w,x\rangle)
:
w \in \mathbb{R}^d,\ \|w\|_0 \le k
\right\},
$$
where
$$
\operatorname{sign}(z)=+1 \text{ for } z\ge 0,
\qquad
\operatorname{sign}(z)=-1 \text{ for } z<0.
$$

\begin{enumerate}
\item Use the Week 4 VC-dimension bound for sparse halfspaces to explain why, if exact empirical risk minimization over $\mathcal{H}_{d,k}$ were computationally available, then the class would be agnostically PAC learnable with sample complexity polynomial in $k$, $\log d$, $1/\varepsilon$, and $\log(1/\delta)$. Then compare this with the brute-force realizable algorithm from Week 4 and explain why the sample-complexity statement and the runtime statement are fundamentally different.

\item Show that deciding empirical agreement for sparse halfspaces is NP-hard when both $d$ and $k$ are part of the input. Reduce from \textsc{Set Cover}. A useful hint is to let coordinates correspond to sets, let negative coordinates of $w$ encode selected sets, let positive examples charge the number of selected sets, and let repeated negative examples reward covered universe elements. Then use the Week 5 learner-to-agreement theorem to conclude that, assuming $\mathrm{NP}\neq\mathrm{RP}$, the family $\mathcal{H}_{d,k}$ is not efficiently properly agnostically PAC learnable when $k$ is part of the input.

\item Run a small experiment comparing:
\begin{itemize}
\item a proper sparse $0/1$ search baseline on very small $d$ and $k$,
\item a convex surrogate method such as hinge-loss minimization with an $\ell_1$ constraint.
\end{itemize}
You may choose the synthetic data yourself. The goal is not a large benchmark; it is to illustrate the difference between exact empirical agreement over sparse classifiers and tractable convex optimization over a larger linear class. Include at least one runtime comparison as $d$ or $k$ increases and one comparison of empirical $0/1$ agreement against the surrogate objective or validation error. State clearly whether your sparse baseline is exact or approximate, and explain whether the surrogate method returns a proper $k$-sparse classifier.
\end{enumerate}

\section{Part B: Surrogate Losses and the Fixed-Feature Barrier (45 points)}

\begin{enumerate}
\item Let $S=((x_1,y_1),\dots,(x_m,y_m))$ with $x_i\in\mathbb{R}^d$ and $y_i\in\{-1,+1\}$. Consider linear predictors $f_w(x)=\langle w,x\rangle$ with the hinge loss
$$
\ell_{\mathrm{hinge}}(f_w(x),y)=(1-yf_w(x))_+,
$$
under an $\ell_1$ constraint
$$
\|w\|_1 \le B.
$$
Write the empirical hinge-loss minimization problem as an explicit linear program by introducing slack variables and the standard positive/negative decomposition of $w$. Then explain carefully why this tractable surrogate optimization problem is \emph{not} the same as solving empirical $0/1$ agreement over $\mathcal{H}_{d,k}$.

\item Consider the one-dimensional domain $\mathcal{X}=\{-M,1\}$ with labels always $+1$. Let
$$
\mathbb{P}[x=1]=1-p,
\qquad
\mathbb{P}[x=-M]=p,
$$
where
$$
0<p<\frac12,
\qquad
M>\frac{1-p}{p}.
$$
For predictors $f_w(x)=wx$, prove the following:
\begin{enumerate}
\item The best $0/1$ classifier has population error at most $p$.
\item Every minimizer of the population hinge risk has $0/1$ error $1-p$.
\end{enumerate}
Conclude that for every $\varepsilon>0$ and every $\alpha<1$ there exists a distribution of this form such that
$$
\inf_w L^{0/1}_{\mathcal{D}}(f_w)\le \varepsilon,
$$
but every hinge-risk minimizer has $0/1$ error greater than $\alpha$.

\item Let
$$
\mathcal{X}=\{-1,+1\}^d,
$$
and for each $I\subseteq[d]$, define the parity
$$
\chi_I(x)=\prod_{i\in I}x_i.
$$
Suppose a fixed feature map
$$
\phi:\{-1,+1\}^d \to \mathbb{R}^D
$$
has the property that for every parity $\chi_I$ there exists some $w_I\in\mathbb{R}^D$ with
$$
\chi_I(x)=\langle w_I,\phi(x)\rangle
\qquad
\text{for all } x\in\{-1,+1\}^d.
$$
Prove that necessarily
$$
D\ge 2^d.
$$
Then give a matching construction showing that
$$
D=2^d
$$
is sufficient. Interpret the result in the context of the Week 5 ``fixed features versus learned features'' discussion.

\item Compare the linear predictor class
$$
\mathcal{F}_{\mathrm{lin}}
=
\left\{
x \mapsto \langle \beta,x\rangle : \beta\in\mathbb{R}^d
\right\}
$$
with the one-hidden-unit two-layer class
$$
\mathcal{F}_{\mathrm{net}}
=
\left\{
x \mapsto v\langle u,x\rangle : u\in\mathbb{R}^d,\ v\in\mathbb{R}
\right\}.
$$
Show that these two classes represent exactly the same set of linear predictors. Then consider squared loss on a fixed sample:
$$
L_{\mathrm{lin}}(\beta)=\frac1m \sum_{i=1}^m (\langle \beta,x_i\rangle-y_i)^2,
$$
$$
L_{\mathrm{net}}(u,v)=\frac1m \sum_{i=1}^m (v\langle u,x_i\rangle-y_i)^2.
$$
Explain why $L_{\mathrm{lin}}$ is convex in $\beta$, why $L_{\mathrm{net}}$ is generally nonconvex in $(u,v)$, and why every global minimizer of the linear problem can be factorized into a global minimizer of the two-layer problem. What does this say about the difference between representational power and optimization geometry?
\end{enumerate}

\section{Part C: AI Usage Report (15 points)}

Write a short report describing how you used AI in this assignment. Explain:
\begin{enumerate}
\item which parts of the assignment used AI,
\item concrete suggestions you accepted,
\item concrete suggestions you rejected or revised,
\item how you verified the submitted proofs and experiment,
\item the five most recent workflow updates you made for this repo and why.
\end{enumerate}

\end{document}
